% O que o ARIMA prevê e o que acontece: a previsão de cada pregão contra o
% retorno realizado (traço da refação de 05/10/2026, jul/2019 a jun/2021).
% A previsão é a mesma com e sem GARCH.
\documentclass[border=6pt]{standalone}
% Preâmbulo comum dos gráficos que explicam os métodos (30/10/2026).
% Dados e números: dados/resultados/metodos_ilustrados_2026-10-30/, gerados por
% ferramentas/MetodosIlustrados.mq5 com as funções do próprio robô. Nenhum
% número é digitado aqui: os do texto vêm de valores.tex (\vX para o texto,
% \pX para coordenada).
\usepackage[T1]{fontenc}
\usepackage{amsmath}
\usepackage{mathptmx}  % Times, como as figuras do TG
\usepackage{pgfplots}
\usepgfplotslibrary{groupplots,dateplot,colormaps,fillbetween}
\usetikzlibrary{arrows.meta,positioning}
\pgfplotsset{compat=1.18}
\definecolor{cmod}{RGB}{0,90,181}     % modelo / ajuste
\definecolor{cdes}{RGB}{220,50,32}    % destaque, limite, valor crítico
\definecolor{cdad}{RGB}{120,120,120}  % dado observado
\newcommand{\arq}[1]{../../dados/resultados/metodos_ilustrados_2026-10-30/#1}
\input{\arq{valores.tex}}
% Figuras 08 a 14 (o robô operando): traço por barra da refação de 05/10/2026,
% resumido por ferramentas/scripts/robo_em_acao.sh. Macros \vR e \pR.
\newcommand{\arqr}[1]{../../dados/resultados/robo_em_acao_2026-10-30/#1}
\input{\arqr{valores.tex}}
\definecolor{cari}{RGB}{230,140,0}   % ARIMA puro
\definecolor{cok}{RGB}{0,140,70}     % compra liberada
\usetikzlibrary{shapes.misc}
\newcommand{\passo}[1]{\tikz[baseline=(n.base)]\node[circle,fill=black!75,text=white,inner sep=1pt,font=\footnotesize\bfseries](n){#1};}
% eixo de datas comum às figuras do traço (jul/2019 a jun/2021)
\pgfplotsset{traco/.style={date coordinates in=x, xmin=2019-07-01, xmax=2021-06-30,
    xtick={2019-07-01,2020-01-01,2020-07-01,2021-01-01,2021-06-30},
    xticklabels={jul/2019,jan/2020,jul/2020,jan/2021,jun/2021}}}
\pgfplotsset{
  base/.style={table/col sep=comma, grid=major, grid style={gray!20},
    tick label style={font=\small}, label style={font=\small},
    title style={font=\small\bfseries, align=center},
    every axis plot/.append style={line width=0.9pt},
    /pgf/number format/use comma, /pgf/number format/1000 sep={.}},
}
% faixa vertical nos pregões do traço em que a condição #1 vale, na cor #2,
% do piso #3 ao teto #4 do eixo
\newcommand{\faixa}[4]{%
  \addplot[name path=topo, draw=none, const plot, forget plot] table[x=data, y expr={(#1) ? #4 : #3}]{\arqr{dia_a_dia.csv}};
  \addplot[name path=piso, draw=none, forget plot] coordinates {(2019-07-01,#3) (2021-06-30,#3)};
  \addplot[#2, forget plot] fill between[of=topo and piso];}

\begin{document}
\begin{tikzpicture}
\begin{axis}[base, traco, name=serie, width=11cm, height=6.4cm, ymin=-12, ymax=12, clip=true,
  ylabel={Retorno do dia (\%)},
  title={\passo{1} No tempo: a previsão (azul) é pequena perto do que acontece (cinza)},
  legend style={at={(0.5,0.02)}, anchor=south, font=\footnotesize, legend columns=2, fill=white, fill opacity=0.9}]
\addplot[ycomb, cdad!70, line width=0.5pt, mark=none] table[x=data, y=real]{\arqr{dia_a_dia.csv}};
\addlegendentry{retorno realizado}
\addplot[cmod, line width=1pt] table[x=data, y=prev]{\arqr{dia_a_dia.csv}};
\addlegendentry{retorno previsto pelo ARIMA}
\node[font=\footnotesize, align=left, anchor=north west, fill=white, fill opacity=0.85, text opacity=1]
  at (rel axis cs:0.01,0.99) {desvio-padrão: previsto $\vRPrevDP\%$, realizado $\vRRealDP\%$\\(dias além de $\pm 12\%$ cortados)};
\end{axis}
\begin{axis}[base, name=disp, at={(serie.east)}, anchor=west, xshift=1.9cm, width=7.6cm, height=7.6cm,
  xmin=-12, xmax=12, ymin=-12, ymax=12, axis equal image, clip=true,
  xlabel={Retorno previsto (\%)}, ylabel={Retorno realizado (\%)},
  title={\passo{2} Um ponto por pregão:\\acertou a direção em \vRPrevAcertos{} de \vRPrevN{} ($\vRPrevTaxa\%$)}]
\fill[cok!10] (axis cs:0,0) rectangle (axis cs:12,12);
\fill[cok!10] (axis cs:0,0) rectangle (axis cs:-12,-12);
\fill[cdes!8] (axis cs:0,0) rectangle (axis cs:-12,12);
\fill[cdes!8] (axis cs:0,0) rectangle (axis cs:12,-12);
\draw[black!40, dashed] (axis cs:-12,-12) -- (axis cs:12,12);
\addplot[only marks, mark=*, mark size=1.1pt, cok, opacity=0.7] table[x=prev, y=real]{\arqr{prev_acerto.csv}};
\addplot[only marks, mark=*, mark size=1.1pt, cdes, opacity=0.7] table[x=prev, y=real]{\arqr{prev_erro.csv}};
\node[cok!70!black, font=\footnotesize\bfseries] at (axis cs:6,10.5) {acerto};
\node[cok!70!black, font=\footnotesize\bfseries] at (axis cs:-6,-10.5) {acerto};
\node[cdes, font=\footnotesize\bfseries] at (axis cs:-6,10.5) {erro};
\node[cdes, font=\footnotesize\bfseries] at (axis cs:6,-10.5) {erro};
\node[black!55, font=\scriptsize, rotate=45, anchor=south] at (axis cs:7,7) {previsão perfeita};
\end{axis}
\node[font=\footnotesize, text=black!75, align=left, anchor=north west, text width=19cm]
  at ([yshift=-0.15cm]serie.outer west |- disp.outer south) {Como ler: se o ARIMA previsse bem, os pontos da direita ficariam sobre a diagonal. Eles formam uma faixa
  vertical estreita: a previsão varia pouco ($\vRPrevDP\%$ de desvio-padrão) e o mercado varia cerca de $\vRRazaoDP$ vezes mais.
  O robô usa só o sinal da previsão (compra se positiva), e a direção acertou em $\vRPrevTaxa\%$ dos pregões, praticamente o mesmo
  que uma moeda. A previsão saiu positiva em $\vRPrevPosPct\%$ dos dias. Os saltos de abril de 2020 são do trecho em que o ADF escolheu $d = 0$ (figura 10).};
\end{tikzpicture}
\end{document}
